\Needspace{12\baselineskip}
\begin{center}
\textit{Besondere Anwendungen auf die Zerlegung von Zahlen in zwei Quadrate, in ein Quadrat und das Doppelte eines Quadrats, in ein Quadrat und das Dreifache eines Quadrats.}
\end{center}

\Needspace{8\baselineskip}
\Article{182}
Wir wenden uns einigen besonderen Fällen zu, sowohl wegen ihrer bemerkenswerten Eleganz als auch wegen der anhaltenden Arbeit, die der berühmte Euler ihnen gewidmet hat. Dadurch haben sie gewissermaßen klassischen Rang erlangt.

I. Durch die Form $xx+yy$ lässt sich eine Zahl nur dann so darstellen, dass $x$ zu $y$ teilerfremd ist (oder in
\SourceLeaf{165}
zwei teilerfremde Quadrate zerlegen), wenn $-1$ ein quadratischer Rest dieser Zahl ist. Umgekehrt ist dies bei allen solchen Zahlen möglich, wenn man sie positiv nimmt. Sei $M$ eine solche Zahl, und seien alle Werte des Ausdrucks $\surd-1\ (\mathrm{mod.}\ M)$ die folgenden: $N,-N,N',-N',N'',-N''$ usw. Dann ist nach Artikel 176 die Form $(M,N,\frac{NN+1}{M})$ der Form $(1,0,1)$ eigentlich äquivalent. Eine eigentliche Transformation der zweiten in die erste sei $x=\alpha x'+\beta y'$, $y=\gamma x'+\delta y'$. Dann sind die Darstellungen der Zahl $M$ durch die Form $xx+yy$, die zu $N$ gehören, diese vier\HistoricalNote{*)}{Dieser Fall ist nämlich offenbar unter (2) von Artikel {180} enthalten.}: $x=\pm\alpha,y=\pm\gamma$; $x=\mp\gamma,y=\pm\alpha$.

Da die Form $(1,0,1)$ ambig ist, ist offenbar auch die Form $(M,-N,\frac{NN+1}{M})$ ihr eigentlich äquivalent, und die erste geht eigentlich in die zweite über, wenn man $x=\alpha x'-\beta y'$, $y=-\gamma x'+\delta y'$ setzt. Daraus ergeben sich vier Darstellungen von $M$, die zu $-N$ gehören: $x=\pm\alpha,y=\mp\gamma$; $x=\pm\gamma,y=\pm\alpha$. Es gibt also acht Darstellungen von $M$, von denen die eine Hälfte zu $N$, die andere zu $-N$ gehört. Sie alle ergeben aber nur \textit{eine} Zerlegung der Zahl $M$ in zwei Quadrate, $M=\alpha\alpha+\gamma\gamma$, wenn wir nur die Quadrate selbst betrachten, ohne ihre Reihenfolge oder die Vorzeichen ihrer Wurzeln zu berücksichtigen.

Gibt es also keine anderen Werte des Ausdrucks $\surd-1\ (\mathrm{mod.}\ M)$ als $N$ und $-N$, was \textit{zum Beispiel} eintritt, wenn $M$ eine Primzahl ist, so lässt sich $M$ nur auf eine Art in zwei teilerfremde Quadrate zerlegen. Da nun $-1$ ein quadratischer Rest jeder Primzahl der Form $4n+1$ ist (Artikel 108) und sich eine Primzahl offenbar nicht in zwei Quadrate zerlegen lässt, die nicht teilerfremd sind, haben wir den Satz:

\textit{Jede Primzahl der Form $4n+1$ lässt sich in zwei Quadrate zerlegen, und zwar nur auf eine Art.}
\[
\begin{gathered}
1=0+1,\quad5=1+4,\quad13=4+9,\quad17=1+16,\\
29=4+25,\quad37=1+36,\quad41=16+25,\quad53=4+49,\\
61=25+36,\quad73=9+64,\quad89=25+64,\quad97=16+81\quad\text{usw.}
\end{gathered}
\]
Diesen sehr eleganten Satz kannte schon Fermat, doch bewies ihn zuerst der berühmte Euler, \textit{Comm. nov. Petr. Bd.}\ V, \textit{für die Jahre} 1754, 1755, \textit{S.}\ 3 ff. In \textit{Bd.}\ IV steht eine Abhandlung zum selben Thema, \textit{S.}\ 3 ff.; damals hatte er die Sache aber noch nicht vollständig abgeschlossen. Siehe besonders Artikel 27.

\SourceLeaf{166}
Lässt sich also eine Zahl der Form $4n+1$ entweder auf mehrere Arten oder gar nicht in zwei Quadrate zerlegen, so ist sie sicher keine Primzahl.

Hat umgekehrt der Ausdruck $\surd-1\ (\mathrm{mod.}\ M)$ außer $N$ und $-N$ noch andere Werte, so gibt es weitere Darstellungen von $M$, die zu diesen gehören. In diesem Fall lässt sich also $M$ auf mehrere Arten in zwei Quadrate zerlegen, \textit{zum Beispiel} $65=1+64=16+49$, $221=25+196=100+121$.

Die übrigen Darstellungen, bei denen $x,y$ keine teilerfremden Werte haben, lassen sich mit unserer allgemeinen Methode leicht finden. Wir bemerken nur: Enthält eine Zahl Faktoren der Form $4n+3$, die sich durch keine Division durch ein Quadrat entfernen lassen (was eintritt, wenn einer oder mehrere dieser Faktoren mit \textit{ungeradem Exponenten} vorkommen), so lässt sie sich auf keine Art in zwei Quadrate zerlegen\HistoricalNote{*)}{Ist die Zahl $M=2^{\mu}Sa^{\alpha}b^{\beta}c^{{\gamma}}\ldots$, wobei $a,b,c$ usw. verschiedene Primzahlen der Form $4n+1$ sind und $S$ das Produkt aller Primfaktoren von $M$ der Form $4n+3$ ist (jede positive Zahl lässt sich in diese Form bringen, indem man $\mu=0$ setzt, wenn $M$ ungerade ist, und $S=1$, wenn $M$ keine Faktoren der Form $4n+3$ enthält), so lässt sich $M$ auf \textit{keine Art} in zwei Quadrate zerlegen, wenn $S$ kein Quadrat ist. Ist dagegen $S$ ein Quadrat, so gibt es $\tfrac12(\alpha+1)(\beta+1)(\gamma+1)$ usw. Zerlegungen von $M$, wenn mindestens eine der Zahlen $\alpha,\beta,\gamma$ usw. ungerade ist, oder $\tfrac12(\alpha+1)(\beta+1)(\gamma+1)\text{ usw.}+\tfrac12$, wenn alle $\alpha,\beta,\gamma$ usw. gerade sind (wobei nur die Quadrate selbst betrachtet werden). Wer mit kombinatorischen Rechnungen etwas vertraut ist, kann den Beweis dieses Satzes ohne große Mühe aus unserer allgemeinen Theorie herleiten. Wir können uns hier bei ihm ebenso wenig wie bei anderen \textit{besonderen Fällen} aufhalten. Vergleiche Artikel 105.}.

II. Durch die Form $xx+2yy$ lässt sich keine Zahl, für die $-2$ ein Nichtrest ist, so darstellen, dass $x$ zu $y$ teilerfremd ist; bei allen übrigen ist dies möglich. Sei $-2$ ein Rest der Zahl $M$ und $N$ ein Wert des Ausdrucks $\surd-2\ (\mathrm{mod.}\ M)$. Dann sind nach Artikel 176 die Formen $(1,0,2)$, $(M,N,\frac{NN+2}{M})$ eigentlich äquivalent. Die erste gehe durch $x=\alpha x'+\beta y'$, $y=\gamma x'+\delta y'$ eigentlich in die zweite über. Dann ist $x=\alpha,y=\gamma$ eine Darstellung der Zahl $M$, die zu $N$ gehört. Außer dieser und $x=-\alpha,y=-\gamma$ gehören keine weiteren zu $N$ (Artikel 180).

Wie oben erkennt man, dass die Darstellungen $x=\pm\alpha,y=\mp\gamma$ zum Wert $-N$ gehören. Alle vier Darstellungen ergeben nur eine Zerlegung von $M$ in ein Quadrat und das Doppelte eines Quadrats. Gibt es außer $N$ und $-N$ keine weiteren Werte des Ausdrucks $\surd-2\ (\mathrm{mod.}\ M)$, so gibt es keine weiteren Zerlegungen. Mit den Aussagen von Artikel 116 folgt daraus leicht der Satz:

\SourceLeaf{167}
\textit{Jede Primzahl der Form $8n+1$ oder $8n+3$ lässt sich in ein Quadrat und das Doppelte eines Quadrats zerlegen, und zwar nur auf eine Art.}
\[
\begin{gathered}
1=1+0,\quad3=1+2,\quad11=9+2,\quad17=9+8,\\
19=1+18,\quad41=9+32.\quad43=25+18,\quad59=9+50,\\
67=49+18,\quad73=1+72,\quad83=81+2,\\
89=81+8,\quad97=25+72\quad\text{usw.}
\end{gathered}
\]
Auch diesen Satz kannte Fermat, wie mehrere ähnliche; den ersten Beweis gab aber der berühmte La Grange, \textit{Suite des recherches d'Arithmétique}, Nouv. Mém. de l'Ac. de Berlin 1775, S. 323 ff. Vieles zum selben Thema hatte bereits der berühmte Euler geklärt, \textit{Specimen de usu observationum in mathesi pura}, Comm. nov. Petr. Bd. VI, S. 185 ff. Ein vollständiger Beweis des Satzes entzog sich aber immer seinen Bemühungen, S. 220. Vergleiche auch die Abhandlung in Bd. VIII (für die Jahre 1760, 1761), \textit{Supplementum quorundam theorematum arithmeticorum}, gegen Ende.

III. Mit einer ähnlichen Methode beweist man, dass sich jede Zahl, für die $-3$ ein quadratischer Rest ist, entweder durch die Form $xx+3yy$ oder durch diese $2xx+2xy+2yy$ darstellen lässt, und zwar so, dass der Wert von $x$ zum Wert von $y$ teilerfremd ist. Da nun $-3$ ein Rest aller Primzahlen der Form $3n+1$ ist (Artikel 119) und die Form $2xx+2xy+2yy$ offenbar nur \textit{gerade} Zahlen darstellen kann, ergibt sich wie oben der Satz:

\textit{Jede Primzahl der Form $3n+1$ lässt sich in ein Quadrat und das Dreifache eines Quadrats zerlegen, und zwar nur auf eine Art.}
\[
\begin{gathered}
1=1+0,\quad7=4+3,\quad13=1+12,\quad19=16+3,\\
31=4+27,\quad37=25+12,\quad43=16+27,\quad61=49+12,\\
67=64+3,\quad73=25+48\quad\text{usw.}
\end{gathered}
\]
Den Beweis dieses Satzes gab zuerst der berühmte Euler in der eben angeführten Abhandlung, \textit{Comm. nov. Petr. Bd.}\ VIII, \textit{S.}\ 105 ff.

Ebenso könnten wir weitergehen und \textit{zum Beispiel} zeigen, dass sich jede Primzahl der Form $20n+1$, $20n+3$, $20n+7$ oder $20n+9$ (für die nämlich $-5$ ein Rest ist) durch eine der Formen $xx+5yy$, $2xx+2xy+3yy$ darstellen lässt: die Primzahlen der Formen $20n+1$ und $20n+9$ durch die erste, die Primzahlen der Formen $20n+3,20n+7$ durch die zweite. Ebenso lässt sich das Doppelte der Primzahlen der Formen $20n+1,20n+9$ durch die Form $2xx+2xy+3yy$ und das Doppelte der Primzahlen der Formen $20n+3,20n+7$ durch die Form $xx+5yy$ darstellen. Doch kann jeder diese Aussage und unendlich viele weitere besondere Aussagen selbst aus den vorigen Ergebnissen und dem herleiten, was wir weiter unten
\SourceLeaf{168}
behandeln werden. --- Wir gehen also zu den \textit{Formen mit positiver Determinante} über. Da ihre Art ganz verschieden ist, je nachdem, ob die Determinante ein Quadrat ist oder nicht, schließen wir die Formen mit quadratischer Determinante zunächst aus und werden sie später gesondert betrachten.
